\chapter{Потенциал центрированной связности и динамический запрет единственности}

Предыдущая классификация оставила после удаления общего скаляра
четырёхмерное комплексное пространство совместимых связностей. Настоящая
глава выполняет последний разрешённый тест этой ветви: может ли уже
имеющийся следовый функционал выбрать в нём единственную ненулевую
юкавскую орбиту без новых коэффициентов?

\section{Сверка с прежними результатами}

В Томе III выравнивающий функционал
\(\Tr(D_F^2-\mathbf 1)^2\) выбрал орбиту малого двухрёберного оператора, но
оставил внешний физический масштаб. В Томе IV было показано, что повышение
степени спектрального полинома меняет радиус стационарной точки, но не
устраняет структурные нулевые направления. Поэтому здесь недостаточно
предъявить один удачный потенциал: требуется классифицировать весь класс,
разрешённый уже зафиксированными физическими блоками.

Это согласуется с общим языком теории инвариантов: симметрия задаёт
порождающие инварианты, но не коэффициенты произвольного инвариантного
полинома. Квадрат нормы отображения момента также зависит от выбора
центрального уровня. См. \path{arXiv:1903.09357},
\path{arXiv:2010.08294} и \path{arXiv:2102.07727}. В спектральной
Стандартной модели параметры конечного оператора Дирака соответствуют
юкавским данным; см. \path{arXiv:1101.2174}.

\section{Центрированное пространство}

Запишем разность связностей в наблюдаемом углу как
\begin{equation}
 C=\sum_{s=1}^{5}z_sP_s,
 \qquad
 (w_1,\ldots,w_5)=\frac1{15}(6,2,3,3,1).
\end{equation}
Удаление общего скаляра эквивалентно условию
\begin{equation}
 \tau_{15}(C)=\sum_{s=1}^{5}w_sz_s=0.
 \label{eq:v5-centered-connection-constraint}
\end{equation}
Его ядро имеет комплексную размерность четыре. Машинный дефект построенного
базиса не превышает \(2.1\cdot10^{-17}\).

Пять наблюдаемых проекторов являются физически различимыми: они отвечают
блокам \(Q_L,L_L,u_R,d_R,e_R\), имеют разные размерности и разные действия
Стандартной модели. Поэтому текущая симметрия не переставляет их.

\section{Общий чётный потенциал}

Даже если дополнительно потребовать чётность по \(C\) и нечувствительность
к фазам отдельных компонент, разрешены пять квадратичных инвариантов
\begin{equation}
 I_s=\tau_{15}(P_sC^*C)=w_s|z_s|^2.
\end{equation}
После ограничения \eqref{eq:v5-centered-connection-constraint} все пять
остаются линейно независимыми. До четвёртой степени общий функционал имеет
вид
\begin{equation}
 V(C)=\sum_s m_s I_s+\sum_{s\le t}\lambda_{st}I_sI_t+\mathrm{const}.
 \label{eq:v5-general-centered-connection-potential}
\end{equation}
Имеются пять независимых квадратичных и пятнадцать квартичных коэффициентов.
Без дополнительной чётности и фазовой слепоты допустимый класс только
расширяется.

\begin{proposition}[Симметрия не фиксирует потенциал]
Ни алгебра Стандартной модели, ни семейный проектор ранга один, ни
единственный след не связывают коэффициенты \(m_s\) и
\(\lambda_{st}\) в
\eqref{eq:v5-general-centered-connection-potential}. Поэтому общий
симметрийно разрешённый потенциал может выбрать требуемое направление
только посредством новых межблочных данных.
\end{proposition}

\section{Наиболее благоприятное следовое ограничение}

Чтобы не закрывать путь одной лишь общей классификацией, запретим вставки
проекторов вручную и оставим только
\begin{equation}
 K=C^*C,\qquad
 p_1=\tau_{15}(K),\qquad
 p_2=\tau_{15}(K^2).
\end{equation}
Тогда общий чётный функционал до четвёртой степени равен
\begin{equation}
 V_{a,b,c}=a\,p_1+b\,p_1^2+c\,p_2+\mathrm{const}.
 \label{eq:v5-trace-only-centered-potential}
\end{equation}
Три функции \(p_1,p_1^2,p_2\) независимы на центрированном пространстве.
После удаления общей нормировки остаются два свободных отношения
коэффициентов.

Имеющийся положительный родительский квадрат даёт \(V=p_1\), чей
единственный минимум есть
\begin{equation}
 C=0.
\end{equation}
Он подавляет юкавскую связность, а не выводит её.

Радиальный функционал
\begin{equation}
 V_{\mathrm{rad}}=(p_1-1)^2
\end{equation}
имеет минимум на взвешенной сфере в \(\mathbb C^4\). Его гессиан имеет одно
положительное и семь нулевых вещественных направлений. Замена единицы на
масштаб \(M^2\) не снимает вырождение и добавляет внешний масштаб.

\section{Проверка выравнивающего функционала}

Самый сильный старый кандидат имеет вид
\begin{equation}
 V_{\mathrm{flat}}(C)
 =\tau_{15}\!\left(C^*C-\mathbf1\right)^2
 =\sum_s w_s\left(|z_s|^2-1\right)^2.
 \label{eq:v5-centered-flattening-potential}
\end{equation}
Нулевой минимум требует \(|z_s|=1\) и
\(\sum_sw_sz_s=0\). Такой многоугольник существует, поскольку наибольший
вес \(6/15\) не превосходит сумму остальных весов \(9/15\).

Машинный аудит предъявляет два различных нулевых минимума с дефектом
центрирования меньше \(1.2\cdot10^{-9}\). Точный гессиан
\eqref{eq:v5-centered-flattening-potential} на центрированном пространстве
имеет пять положительных и три нулевых собственных значения. Следовательно,
минимумы образуют как минимум трёхмерное вещественное семейство.

\begin{theorem}[Динамическая неединственность]
Текущий положительный родительский квадрат выбирает нулевую связность.
Общий симметрийно разрешённый квартичный потенциал содержит независимые
межблочные коэффициенты. Даже наиболее ограничительный выравнивающий
функционал одного следа имеет непрерывное трёхмерное семейство нулевых
минимумов и требует внешнего единичного масштаба. Поэтому существующая
архитектура не выбирает единственную ненулевую юкавскую орбиту.
\end{theorem}

\section{Почему центральный уровень не является решением}

Можно заменить нулевое условие отображения момента на
\(\mu=\zeta\), где \(\zeta\) лежит в центре пятиблочной алгебры. Но
следовое условие \(\tau_{15}(\zeta)=0\) оставляет четыре вещественных
центральных параметра. Это точно повторяет число относительных блочных
направлений, которое требовалось вывести.

\begin{nogocriterion}[Запрет скрытого уровня]
Нельзя выбирать юкавскую орбиту подстановкой ненулевого центрального уровня,
набора квадратичных масс, квартичных коэффициентов или спектрального
профиля, если эти данные не выведены независимо от требуемых масс и
смешиваний.
\end{nogocriterion}

\section{Итог ветви}

Связывающий родитель Мориты сохраняется как точная кинематическая
архитектура носителя, двух действий, следа и относительной кривизны. Но
происхождение физической юкавской связности внутри этой архитектуры
недоопределено.

\begin{protocol}
\begin{enumerate}
 \item размерность центрированной связности: \textbf{четыре комплексных};
 \item независимые квадратичные блочные инварианты: \textbf{пять};
 \item независимые квартичные произведения: \textbf{пятнадцать};
 \item коэффициенты общего потенциала из текущей симметрии:
 \textbf{не выведены};
 \item положительная родительская норма: \textbf{выбирает нуль};
 \item радиальный следовый потенциал: \textbf{семь нулевых мод};
 \item выравнивающий потенциал: \textbf{три нулевые моды};
 \item ненулевой центральный уровень: \textbf{внешний выбор};
 \item единственная юкавская орбита: \textbf{не получена};
 \item связывающий носитель Мориты: \textbf{сохранён};
 \item физическое замыкание: \textbf{не достигнуто}.
\end{enumerate}
\end{protocol}

Повторное открытие возможно только после независимого вывода новой
симметрии, центрального уровня, спектрального профиля или фермионного
принципа, устраняющего четыре относительных блочных направления.

Машинный сертификат сохранён в
\path{s2t_v5_centered_connection_potential_gate.py} и
\path{s2t_v5_centered_connection_potential_gate_results.json}.